%!TEX root = ../template.tex
%%%%%%%%%%%%%%%%%%%%%%%%%%%%%%%%%%%%%%%%%%%%%%%%%%%%%%%%%%%%%%%%%%%%
%% abstrac-en.tex
%% NOVA thesis document file
%%
%% Abstract in English([^%]*)
%%%%%%%%%%%%%%%%%%%%%%%%%%%%%%%%%%%%%%%%%%%%%%%%%%%%%%%%%%%%%%%%%%%%

\typeout{NT FILE abstrac-en.tex}%
In recent years, the impact of climate change has manifested in more frequent and
severe wildfires, posing significant threats to ecosystems, properties, and human
lives. Effective vegetation management around these
areas is crucial for wildfire risk reduction. Recognizing the fire risk associated with
various vegetation types is essential for planning preventative measures to protect homes
from the devastating impact of wildfires.

This thesis plan aims to find a way to automate the assessment of fire risks using machine learning techniques,
focusing on the vegetation surrounding residential areas.
The classifications provided by \gls{adai},
covers over 2.000 isolated house locations in Continental Portugal, will be used as
ground truth for our study. Aerial images of the houses and the surrounding vegetation
are obtained through the open data API of the \gls{dgt}.

Related works such as the Hydra and Feature Aggregation Convolutional Neural Network (FACNN),
present highly effective model architectures in aerial and satellite image classifications.
These models are based on Convolutional Neural Networks (CNNs) and transfer learning,
providing a solid foundation for our problem.

Additionally, we outline a comprehensive and structured work plan involving data
collection, preprocessing, model implementation, training, and evaluation.
This approach aims to ensure the development of a scalable and
efficient solution for fire risk assessment.

By automating the identification of vegetation risk classes, this research seeks
to reduce the time and effort required for risk assessment, enabling more timely
and effective mitigation strategies. Ultimately, this study contributes to the
broader goal of utilizing advanced technology to combat the growing threat of
wildfires in a changing climate.

% Palavras-chave do resumo em Inglês
% \begin{keywords}
% Keyword 1, Keyword 2, Keyword 3, Keyword 4, Keyword 5, Keyword 6, Keyword 7, Keyword 8, Keyword 9
% \end{keywords}
\keywords{
  Machine Learning \and
  Supervised Learning \and
  Deep Learning \and
  Convolutional Neural Networks \and
  Remote Sensing
}
